\documentclass[11pt]{letter}
\usepackage[margin=1in]{geometry}
\usepackage{hyperref}

\signature{Elliot Tower \\ Independent Researcher \\ \texttt{elliot@elliottower.ai}}
\address{Elliot Tower \\ Independent Researcher \\ elliot@elliottower.ai}

\begin{document}
\begin{letter}{Editorial Office \\ \emph{Genetic Epidemiology} \\ Wiley}

\opening{Dear Editors,}

Please consider our manuscript, ``Abundance or Activity? Mechanism-Dependent Validity of \emph{Cis}-pQTL Mendelian Randomization,'' for publication as an Original Research Article in \emph{Genetic Epidemiology}.

This study addresses a methodological question central to genetic epidemiology: when does a \emph{cis}-pQTL instrument provide a valid proxy for a pharmacologic intervention? Current evaluations of pQTL-based Mendelian randomization for drug-target validation pool all drug mechanisms and report composite performance. This manuscript shows that such pooling obscures a mechanism-dependent boundary of validity: \emph{cis}-pQTL MR is informative when the drug perturbs protein abundance, but structurally uninformative when the drug acts by blocking protein function without changing measured protein level.

The analysis is a pre-registered, outcome-blind evaluation of a frozen zero-parameter classifier across 195 Phase~III drug-target pairs. The main contribution is not simply a benchmark, but an instrument-validity framework that decomposes composite pQTL performance into interpretable mechanistic strata and helps explain why recent pQTL benchmarks have reported weaker enrichment than GWAS-based approaches.

The manuscript is well aligned with \emph{Genetic Epidemiology} because it focuses on the validity, interpretation, and limitations of genetic instruments rather than on a single therapeutic area. All data sources are public, the full protocol history was SHA-256 hashed before execution, and code and classification tables are available in a public repository. The manuscript is not under consideration elsewhere.

\closing{Sincerely,}

\end{letter}
\end{document}
